\input{../article_base.tex}
\title{פתרון מטלה 03 --- תבניות ריבועיות ומספרים P־אדיים, 80507}

\DeclareMathOperator{\rank}{rank}

% chktex-file 44

\begin{document}
\maketitle
\maketitleprint[red]

\question{}
יהי $p \ne 2$ ראשוני ויהי $(V, q)$ מרחב־ריבועי לא מנוון על $\FF_p$, ונניח ש־$\dim V = 2$ וש־$\Bb$ בסיס של $V$. \\
נראה ש־$q(V \setminus \{ 0 \}) = \FF_p^*$ אם ורק אם $-\det{[g_q]}_\Bb \notin {(\FF_p^*)}^2$.
\begin{proof}
	בכיתה ראינו ש־$- \det {[g_q]}_\Bb \in {(\FF_p^*)}^2$ אם ורק אם $(V, q)$ איזוטרופי ולכן מספיק להראות ש־$q(V \setminus \{ 0 \}) = \FF_p^*$ אם ורק אם $(V, q)$ לא איזוטרופי. \\
	נזכיר שהמרחב איזוטרופי אם קיים $0 \ne v \in V$ כך ש־$q(v) = 0$, אם נניח ש־$0 \notin q(V \setminus \{ 0 \})$ בפרט $(V, q)$ לא איזוטרופי מהגדרה. \\
	מהצד השני נניח ש־$(V, q)$ לא איזוטרופי ולכן $q(V \setminus \{ 0 \}) \subseteq \FF_p^*$ אבל הוכחנו בכיתה שגם $\FF_p^* \subseteq q(V \setminus \{ 0 \})$ וסיימנו.
\end{proof}

\question{}
יהי $p \ne 2$ ראשוני ו־$(V, q)$ מרחב־ריבועי לא מנוון מעל $\FF_p$, ונניח גם ש־$\dim V = 3$ וכן ש־$r \in \FF_p^* \setminus {(\FF_p^*)}^2$. \\
נראה שקיים $a \in \{1, r\}$ ובסיס $\Bb$ של $V$ כך שמתקיים,
\[
	{[g_q]}_\Bb
	= \begin{pmatrix}
		0 & \frac{1}{2} & 0 \\
		\frac{1}{2} & 0 & 0 \\
		0 & 0 & a
	\end{pmatrix}
\]
\begin{proof}
	ממשפט קיום תת־מרחב היפרבולי נסמן $U \subseteq V$ תת־מרחב ממימד $2$ כך ש־$(U, q|_U) \cong (\FF_p^2, h)$ ולכן קיים בסיס $\Bb' = (b_1, b_2)$ של $U$ כך שמתקיים,
	\[
		{[g_q|_U]}_{\Bb'}
		= \begin{pmatrix}
			0 & \frac{1}{2} \\
			\frac{1}{2} & 0
		\end{pmatrix}
	\]
	עתה נרחיב את $\Bb'$ ל־$\Bb = (b_1, b_2, b_3)$ כאשר $b_3$ נבחר כך ש־$0 \ne b_3 \notin U$ וכן $q(b_3) \in \{1, r\}$, מובטח לנו שקיים כזה ממשפט מיון תבניות מעל $\FF_p$.
	נקבל אם כך ש־${[g_q]}_\Bb$ הוא בדיוק מהצורה המבוקשת.
\end{proof}

\question{}
נגדיר את התבנית הריבועית $q : \FF_7^3 \to \FF_7$ על־ידי $q(x_1, x_2, x_3) = x_1^2 + x_2^2 + x_3^2$. \\
נמצא $a \in \{1, -1\}$ ו־$\Bb$ כבשאלה הקודמת כך ש־${[g_q]}_\Bb$ מתקבל כפי שהגדרנו בשאלה זו.
\begin{solution}
	נמצא $v \in \FF_7^3$ כך ש־$q(v) = 0$.
	נבחין (מבחינת לוח כפל מתאים) שמתקיים $q(2, 3, 0) = 2^2 + 3^2 = -1$ ולכן בהתאמה נגדיר $v = (2, 3, 1)$ ונקבל $q(v) = 0$.
	נבחין כי $g_q(x, y) = x_1 * y_1 + \cdots + x_3 y_3$ ובפרט $g_q((1, 0, 0), (2, 3, 1)) = 2$ ונגדיר את $w = \frac{1}{4} (1, 0, 0) = (2, 0, 0)$.
	נקבל אם כך,
	\[
		g_q(v, w)
		= \frac{1}{4} \cdot 2
		= \frac{1}{2}
	\]
	נבחין כי גם $q(w) = 0$ ולכן אם נגדיר $\Bb' = (v, w)$ נקבל,
	\[
		{[g_q|_{\Sp \Bb'}]}_{\Bb'}
		= \begin{pmatrix}
			0 & \frac{1}{2} \\
			\frac{1}{2} & 0
		\end{pmatrix}
	\]
	ונשאר למצוא וקטור $u \in V \setminus \Sp \Bb'$ אורתוגונלי לתת־מרחב זה ו־$q(u) = 0$.
	על־ידי בדיקה ישירה נקבל ש־$e_3$ הוא וקטור $e_3 \notin \Sp \Bb'$ ונשאר לבצע חפיפה ל־$\Cc = (b_1, b_2, e_3)$,
	\[
		{[g_q]}_\Bb
		= \begin{pmatrix}
			0 & \frac{1}{2} & 1 \\
			\frac{1}{2} & 0 & 0 \\
			1 & 0 & 1
		\end{pmatrix}
	\]
	ולכן מספיק לקחת את $\Bb = (v, w, e_3 - \frac{1}{2} w)$ וקיבלנו את המבוקש.
\end{solution}

\question{}
יהי $p \ne 2$ ראשוני ו־$(V, q)$ מרחב־ריבועי לא מנוון מעל $\FF_p$ ממימד $8$. \\
נוכיח שקיים תת־מרחב $U$ של $V$ כך ש־$\dim U = 3$ וגם $q|_U = 0$.
\begin{proof}
	נתון ש־$\dim v \ge 3$ ולכן המרחב איזוטרופי ויהי קיים $v_1 \in V$ כך ש־$v_1 \ne 0$ וגם $q(v_1) = 0$.
	נגדיר אם כך $V_0 = {\{ v_0 \}}^\perp$ ולכן $\dim V_0 = \dim V - 1$.
	מההנחה שלנו גם $(V_0, q|_{V_0})$ ממימד גדול מ־$2$ ולכן קיים $v_1 \in V_0$ כך ש־$q|_{V_0}(v_1) = 0$, ונגדיר את $V_1 = {\{ v_1 \}}^\perp$.
	נמשיך כך רקורסיבית עד שנקבל $\dim V_n = 2$ ושם נעצור, ונגדיר את $U = \{ v_0, \ldots, v_{n - 1} \}$, אז מתקיים $q(u) = 0$ לכל $u \in U$ ישירות מהגדרת $v_i$. \\
	נבחין כי הוכחנו כאן שקיים תת־מרחב כזה ממימד $\dim V - 3$ ובפרט קיים עבור $\FF_p^8$ תת־מרחב ממימד $3$ המקיים את הטענה.
\end{proof}

\question{}
יהיו $(V_1, q_1), (V_2, q_2)$ מרחבים־ריבועיים, ונגדיר פונקציה $q : V_1 \times V_2 \to \FF$ על־ידי,
\[
	q(v_1, v_2)
	= q_1(v_1) + q_2(v_2)
\]
נראה ש־$q$ היא תבנית ריבועית.
\begin{proof}
	נסמן $g_1, g_2$ התבניות הבילינאריות הסימטריות המתאימות ל־$q_1, q_2$ בהתאמה, ונגדיר את $g : {(V_1 \times V_2)}^2 \to \FF$ על־ידי,
	\[
		g((v_1, v_2), (u_1, u_2))
		= g_1(v_1, u_1) + g_2(v_2, u_2)
	\]
	ונראה שזו תבנית בילינארית סימטרית מעל $\FF$.
	מהגדרה ברור כי $g$ סימטרית, ולכן נשאר להראות לינאריות בערך השני.
	\begin{align*}
		g((v_1, v_2), \alpha (u_1, u_2) + \beta (w_1, w_2))
		= & g((v_1, v_2), (\alpha u_1 + \beta w_1, \alpha u_2 + \beta u_2)) \\
		= & g_1(v_1, \alpha u_1 + \beta w_1) + g_2(v_2, \alpha u_2 + \beta w_2) \\
		= & \alpha g_1(v_1, u_1) + \beta g_1(v_1, w_1) + \alpha g_1(v_1, w_1) + \beta g_1(v_1, w_1) \\
		  & + \alpha g_2(v_2, u_2) + \beta g_2(v_2, w_2) + \alpha g_2(v_2, w_2) + \beta g_2(v_2, w_2) \\
		= & \alpha g((v_1, v_2), (u_1, u_2)) + \beta g((v_1, v_2), (w_1, w_2))
	\end{align*}
	כאשר השתמשנו מספר פעמים בלינאריות של $g_1, g_2$ ובהגדרת $g$ המופיעה למעלה.
	מתקיים,
	\[
		g((v_1, v_2), (v_1, v_2))
		= g_1(v_1, v_1) + g_2(v_2, v_2)
		= q_1(v_1) + q_2(v_2)
		= q(v_1, v_2)
	\]
	ומצאנו ש־$g$ היא התבנית הריבועית המתאימה ל־$q$, ומעידה על כך ש־$q$ היא תבנית ריבועית.
\end{proof}
\begin{definition}[מכפלה ישרה של מרחבים־ריבועיים]
	למרחב הריבועי $(V_1 \times V_2, q)$ נקרא המכפלה הישרה של $(V_1, q_1), (V_2, q_2)$ ונסמן $(V_1, q_1) \times (V_2, q_2)$.
\end{definition}

\question{}
יהיו $(V_1, q_1), (V_2, q_2)$ מרחבים־ריבועיים ויהי $(V_1, q_1) \times (V_2, q_2)$ מרחב המכפלה הישרה שלהם.
נגדיר את,
\begin{align*}
	& \pi_1 : (V_1, q_1) \to (V_1, q_1) \times (V_2, q_2), \qquad \forall v \in V_1,\ \pi_1(v) = (v, 0) \\
	& \pi_2 : (V_2, q_2) \to (V_1, q_1) \times (V_2, q_2), \qquad \forall v \in V_2,\ \pi_2(v) = (0, v)
\end{align*}
נראה ש־$\pi_1, \pi_2$ הן איזומטריות וכן שמתקיים $V_1 \times V_2 = \pi_1(V_1) \overset{\perp}{\oplus} \pi_2(V_2)$.
\begin{proof}
	\[
		q_1(v)
		= q_1(v) + 0
		= q_1(v) + q_2(0)
		= q(v, 0)
	\]
	אבל מהגדרה $\pi_1(v) = (v, 0)$ ולכן $q_1(v) = q(\pi_1(v))$ וקיבלנו כי $\pi_1$ איזומטריה, ונקבל בטיעון זהה שגם $\pi_2$ איזומטריה.

	ברור כי $V_1 \times V_2 = V_1 \times \{ 0 \} \oplus \{ 0 \} \times V_2$ ולכן $V_1 \times V_2 = \pi_1(V_1) \times \pi_2(V_2)$, ונשאר להראות שהמרחבים גם אורתוגונליים.
	יהי $v_1 \in \pi_1(V_1), v_2 \in \pi_2(V_2)$, אז מתקיים,
	\[
		g_q(v_1, v_2)
		= \frac{1}{2} (q(v_1 + v_2) - q(v_1) - q(v_2))
		= \frac{1}{2} (q(v_1', v_2') - q_1(v_1') - q_2(v_2'))
		= \frac{1}{2} (q_1(v_1') + q_2(v_2') - q_1(v_1') - q_2(v_2'))
		= 0
	\]
	כאשר $v_1 = (v_1', 0), v_2 = (0, v_2')$, ותוך שימוש בהגדרות, ומצאנו שאכן $V_1 \times V_2 = \pi_1(V_1) \overset{\perp}{\oplus} \pi_2(V_2)$.
\end{proof}

\question{}
יהי $p \ne 2$ ראשוני כך ש־$-1 \in {(\FF_p^*)}^2$ ויהי $r \in \FF_p^* \setminus {(\FF_p^*)}^2$.
נראה כי לכל $k, m, n \in \NN \cup \{ 0 \}$ המרחבים המטריים,
\[
	(\FF_p, d_1) \times {(\FF_p^2, h)}^m,
	\qquad
	(\FF_p^k, d_{r, \ldots, r}) \times {(\FF_p^2, h)}^n
\]
אינם איזומורפיים.
\begin{proof}
	נבחין כי מהגדרה אם $1 + 2m \ne k + 2n$ אז אף איזומטריה לא יכולה להיות חד־חד ערכית ועל ולכן סיימנו, ונניח ש־$1 + 2m = k + 2n$ בלבד.
	על־ידי שימוש בשאלה הקודמת נוכל להסיק שמרחבי מכפלה איזומורפיים אם ורק אם ההטלות איזומורפיות, ולכן מספיק שנראה שהטלות של שני המרחבים הן לא איזומורפיות.
	אם $k = 1$ אז נקבל שהמרחב איזומורפי אם ורק אם $(\FF_p, d_1) \simeq (\FF_p, d_r)$, אבל אנו כבר יודעים כי הם לא איזומורפיים ולכן סיימנו. \\
	נותר לבחון את המקרה $k > 1$, בפרט $n < m$ ונוכל לבחון רק את ההטלה של שני המרחבים ל־$n$ המרחבים הראשונים, כלומר נוכיח שאין איזומורפיזם של המרחבים,
	\[
		(\FF_p, d_1) \times {(\FF_p^2, h)}^{m - n},
		\qquad
		(\FF_p^k, d_{r, \ldots, r})
	\]
	בפרט נקבל $k = 2(m - n) + 1$, כלומר אם $k$ זוגי אז סיימנו, אחרת נוכל ממשפט שקילות למכפלת מרחבים היפרבוליים לצמצם עוד את המרחבים שאנו בוחנים למרחבים ואף להניח $k = 3$ בלבד,
	\[
		(\FF_p, d_1) \times (\FF_p^2, h)
		\qquad
		(\FF_p^k, d_{r, \ldots, r})
		\cong (\FF_p, d_r) \times (\FF_p^2, h)
	\]
	ולכן סיימנו.
\end{proof}

\question{}
יהי $\FF$ שדה סופי ותהי $\lVert \cdot \rVert : \FF \to \RR^+$ נורמה מעליו.
נראה ש־$\lVert \cdot \rVert$ טריוויאלית.
\begin{proof}
	יהי $a \in \FF$, אם $a = 0$ אז $\lVert a \rVert = 0$ מהגדרה.
	נניח ש־$a \ne 0$, אז לכל $n \in \NN$ מתקיים ${\lVert a \rVert}^n = \lVert a^n \rVert$.
	ידוע כי $\FF$ שדה סופי ולכן קיים $p$ ראשוני כך ש־$\FF = \FF_{p^m}$ לאיזשהו $m \in \NN$, ולכן,
	\[
		{\lVert a \rVert}^{p^m}
		= \lVert a^{p^m} \rVert
		= \lVert 1 \rVert
		= 1
	\]
	כאשר $a^{p^m} = a$ כתוצאה ידועה של תורת השדות.
	בפרט $\lVert a \rVert = 1$ בדיוק בעזרת שורש, וקיבלנו שהנורמה היא טריוויאלית.
\end{proof}

\question{}
יהי $\FF \in \{ \QQ, \RR, \CC \}$ ויהי $\alpha \in \RR$ כך ש־$0 < \alpha < 1$.
נגדיר את,
\[
	\lVert \cdot \rVert : \FF \to \RR^+,
	\qquad
	\lVert a \rVert = {| a |}^\alpha
\]
נראה כי $\lVert \cdot \rVert$ היא נורמה ארכימדית.
\begin{proof}
	משקילות מספיק למצוא $n \in \NN \subseteq \FF$ כך שמתקיים $\lVert n \rVert > 1$, בשלושת השדות נקבל גם $| n | = n$ ולכן מספיק למצוא,
	\[
		n^\alpha > 1
	\]
	נזכור כי העלאה בחזקה גדולה מ־$1$ משמרת סדר עבור טבעיים גדולים מ־$1$ ולכן נקבל גם,
	\[
		n^\alpha > 1
		\iff n > n^{\frac{1}{\alpha}}
		= 1
	\]
	אבל $n = 2$ מקיים בדיוק טענה זו בכל השדות הנתונים.
\end{proof}

\question{}
יהי $(\FF, \lVert \cdot \rVert)$ שדה נורמי ויהי $a \in \FF$. \\
נראה ש־$\lVert a \rVert < 1$ אם ורק אם $\lim_{n \to \infty} a^n = 0$.
\begin{proof}
	נניח ש־$\lVert a \rVert < 1$.
	נסמן $\alpha \in \RR^+$ כך ש־$\lVert a \rVert < \alpha < 1$.
	בהתאם באינדוקציה נקבל גם $0 \le {\lVert a \rVert}^n < \alpha^n$ ומהרחבה של משפט הסנדוויץ' לגבולות על מרחבים מטריים נקבל $\lim_{n \to \infty} a^n = 0$ בדיוק.

	נניח ש־$\lim_{n \to \infty} a^n = 0$, מהגדרה לכל $\varepsilon > 0$ קיים $M \in \NN$ כך שלכל $n > M$ מתקיים,
	\[
		\lVert a^n - 0 \rVert
		< \varepsilon
	\]
	ומתכונות נורמה נקבל שגם $\lVert a^n \rVert = {\lVert a \rVert}^n < \varepsilon$.
	נבחר $\varepsilon = 1$ וכן $n = M + 1$ המתאים לו ונקבל $\lVert a \rVert < 1^{\frac{1}{M + 1}} = 1$, כלומר $\lVert a \rVert < 1$.
\end{proof}

\question{}
יהיו $(\FF, \lVert \cdot \rVert)$ שדה נורמי לא ארכימדי, $a \in \FF$, $r \in \RR, r > 0$ וכן $b \in B(a, r)$ עבור,
\[
	B(a, r)
	= \{ x \in \FF \mid \lVert x - a \rVert < r \}
\]
נוכיח כי $B(a, r) = B(b, r)$.
\begin{proof}
	נתון כי $\lVert b - a \rVert < r$ ולכן מסימטריה גם $a \in B(b, r)$ ובפרט מספיק להוכיח את ההכלה $B(a, r) \subseteq B(b, r)$ בלבד (על־ידי החלפת תפקידים). \\
	יהי $x \in B(a, r)$, כלומר $\lVert x - a \rVert < r$, וראינו כי עבור שלושת הערכים,
	\[
		\lVert x - a \rVert < r,
		\quad
		\lVert x - b \rVert,
		\quad
		\lVert a - b \rVert < r
	\]
	יש שני ערכים שווים, ראינו כי שניים מהם חסומים על־ידי $r$ ולכן מעיקרון שובך היונים גם $\lVert x - b \rVert < r$ ובהתאם $x \in B(b, r)$ ומצאנו הכלה כפי שרצינו.
\end{proof}

\question{}
יהיו $(\FF, \lVert \cdot \rVert)$ שדה נורמי לא ארכימדי, $a \in \FF, r \in \RR, 0 < r$.

\subquestion{}
נראה שהקבוצה $U = \{ x \in \FF \mid \lVert x - a \rVert \le r \}$ פתוחה.
\begin{proof}
	במרחב מטרי נגדיר קבוצה כפתוחה אם לכל נקודה שלה יש כדור פתוח מוכל בה, נראה זאת.
	יהי $x \in U$ ונגדיר את הקבוצה $B = B(x, r)$. \\
	אם $x \in B(a, r)$, כלומר $\lVert x - a \rVert < r$ אז נובע מהשאלה הקודמת ש־$B \subseteq U$. \\
	נניח ש־$\lVert x - a \rVert = r$ ויהי $y \in B(x, r)$, אז מתקיים,
	\[
		\lVert y - a \rVert
		= \lVert y - x + x - a \rVert
		\le \max\{ \lVert y - x \rVert, \lVert a - x \rVert \}
		= r
	\]
	ולכן $B(x, r) \subseteq U$ וקיבלנו שאכן היא קבוצה פתוחה.
\end{proof}

\subquestion{}
נראה שהקבוצה $C = B(a, r)$ סגורה.
\begin{proof}
	נראה שאם $x \in \FF \setminus C$ אז קיים כדור פתוח $x \in B \subseteq C^C$.
	נגדיר את $B = B(x, r)$, אז $x \in B$ ונשאר להראות שהוא זר ל־$C$.
	נניח בשלילה שקיים $y \in C \cap B$, כלומר $\lVert y - a \rVert < r, \lVert y - x \rVert < r$, אז נקבל ש־$\lVert x - a \rVert$ שווה לאחד המרקים ובפרט $< r$, כלומר $x \in C$ בסתירה לבחירה שלנו.
	נסיק שאכן $B \subseteq C^C$ ובהתאם $C$ סגורה.
\end{proof}

\question{}
יהי $(\FF, \lVert \cdot \rVert)$ שדה נורמי לא ארכימדי ויהיו $a, b \in \FF$ כך ש־$a \ne b$. \\
נראה שקיימות קבוצות פתוחות $S, T \subseteq \FF$ כך שמתקיים,
\[
	S \cup T = \FF,
	\quad
	S \cap T = \emptyset,
	\quad
	a \in S,
	\quad
	b \in T
\]
\begin{proof}
	נסמן $r = \lVert a - b \rVert$, ולכן $r > 0$ בלבד.
	נגדיר את $S = B(a, r)$ וכן $T = S^C$, לכן $S \cup T = \FF$ וכן $S \cap T = \emptyset$. \\
	ברור לנו כי $a \in S$, וכן $\lVert a - b \rVert \not< r$ ולכן גם $b \in T$. \\
	לבסוף $S$ פתוחה מהגדרתה ככדור פתוח, אבל מהשאלה הקודמת נקבל שהיא גם סגורה ולכן $T$ פתוחה אף היא.
\end{proof}

\end{document}
